% Copyright 2025 Diego Correa Tristain <algoritmia@labormedia.cl>
%
% Permission is hereby granted, free of charge, to any person obtaining a
% copy of this document and associated documentation files (the "Document"),
% to deal in the Document without restriction, including without limitation
% the rights to use, copy, modify, merge, publish, distribute, sublicense,
% and/or sell copies of the Document, and to permit persons to whom the
% Software is furnished to do so, subject to the following conditions:

% The above copyright notice and this permission notice shall be included in
% all copies or substantial portions of the Document.
%
% THE DOCUMENT IS PROVIDED "AS IS", WITHOUT WARRANTY OF ANY KIND, EXPRESS
% OR IMPLIED, INCLUDING BUT NOT LIMITED TO THE WARRANTIES OF MERCHANTABILITY,
% FITNESS FOR A PARTICULAR PURPOSE AND NONINFRINGEMENT. IN NO EVENT SHALL THE
% AUTHORS OR COPYRIGHT HOLDERS BE LIABLE FOR ANY CLAIM, DAMAGES OR OTHER
% LIABILITY, WHETHER IN AN ACTION OF CONTRACT, TORT OR OTHERWISE, ARISING
% FROM, OUT OF OR IN CONNECTION WITH THE SOFTWARE OR THE USE OR OTHER
% DEALINGS IN THE SOFTWARE.

\documentclass[11pt]{article}

\usepackage{amsmath, amssymb, amsfonts}
\usepackage{mathtools}
\usepackage{hyperref}
\usepackage{fullpage}
\usepackage{bm}
\usepackage{enumitem}
\usepackage{mathrsfs}
\usepackage[backend=biber,style=authoryear]{biblatex} % Load biblatex

\addbibresource{references.bib}

\title{\textbf{Reaction--Diffusion Exchange Economies with Aggregated\\
Cobb--Douglas Preferences for \(\mathbf{n \ge 2}\) Goods}}
\author{Diego Correa Tristain \\ \texttt{algoritmia@labormedia.cl}}

\begin{document}

\maketitle

\begin{abstract}
This document extends the two-good, two-stage convex zero-sum exchange model introduced in
\emph{Emergent Properties of Distributed Agents with Two-Stage Convex Zero-Sum Optimal Exchange Network} \parencite{Tristain2024}
to an economy with \(n \ge 2\) goods. The extension is made possible by aggregating
\emph{pairwise} Cobb--Douglas preference parameters into a globally consistent
Cobb--Douglas exponent vector \(\bm{\beta}^{(i)}\) for each agent \(i\).
The aggregation is exact when a cycle-consistency condition holds; otherwise it can be posed as a least-squares
projection in log-ratio space. The resulting \(n\)-good exchange economy preserves convexity and tractability, and
supports reaction--diffusion dynamics in the sense of Turing-type pattern formation \parencite{Turing1952}.
\end{abstract}

\section{Introduction}

We consider an economy with:
\begin{itemize}[leftmargin=*]
  \item \(m \ge 2\) agents indexed by \(i \in \{1,\dots,m\}\),
  \item \(n \ge 2\) goods (``species'') indexed by \(a \in \{1,\dots,n\}\),
  \item discrete time (rounds) indexed by \(t \in \{0,1,2,\dots\}\).
\end{itemize}

Each agent \(i\) holds an endowment vector
\[
\bm{e}_i^t = (e_{i,1}^t,\dots,e_{i,n}^t)^\top \in \mathbb{R}_+^n,
\]
and the system satisfies a global conservation law across agents in the \emph{diffusion (exchange) phase}:
\[
\sum_{i=1}^m \bm{e}_i^t = \bm{E} \in \mathbb{R}_+^n \qquad \forall t.
\]
In the original two-good model, bilateral exchanges were written as a zero-sum constraint at the edge level.
For \(n\) goods, a bilateral exchange between agents \(i\) and \(j\) can be expressed as an exchange vector
\(\bm{\varepsilon}_{i \to j}^{t} \in \mathbb{R}^n\) such that
\[
\bm{\varepsilon}_{i \to j}^{t} + \bm{\varepsilon}_{j \to i}^{t} = \bm{0},
\]
ensuring that trades only redistribute goods among agents.

The dynamics have two phases each round:
\begin{enumerate}[leftmargin=*]
  \item \textbf{Reaction:} each agent applies endogenous transformation rules (technologies) to its own endowment,
        mapping some goods into others under feasibility constraints.
  \item \textbf{Diffusion:} agents exchange goods via local, utility-maximizing trades subject to budgets, producing
        a redistribution consistent with conservation.
\end{enumerate}
This mirrors reaction--diffusion systems in which reaction reshapes local composition while diffusion transports
quantities across the network \parencite{Turing1952}. The key additional ingredient in the economic setting is that
diffusion is not random transport but \emph{Pareto-improving exchange} governed by convex preferences.

\section{Pairwise Cobb--Douglas Preferences}

Fix an agent \(i\). For each ordered pair of \emph{distinct} goods \((a,b)\) the agent specifies a pairwise
Cobb--Douglas utility over the two-dimensional projection \((x_a,x_b)\):
\[
u_{i,(a,b)}(x_a,x_b) = x_a^{\alpha^{(i)}_{ab}} x_b^{1-\alpha^{(i)}_{ab}},
\qquad \alpha^{(i)}_{ab} \in (0,1).
\]
A natural symmetry requirement is
\[
\alpha^{(i)}_{ba} = 1 - \alpha^{(i)}_{ab},
\]
so that the same indifference curves are obtained regardless of the direction in which the pair is written.

\subsection{Preference ratios}

Define the pairwise \emph{intensity ratio} (good \(b\) relative to \(a\)):
\[
r^{(i)}_{ab} \coloneqq \frac{1-\alpha^{(i)}_{ab}}{\alpha^{(i)}_{ab}} > 0.
\]
In a global Cobb--Douglas representation, these ratios correspond to exponent ratios:
\[
r^{(i)}_{ab} = \frac{\beta_b^{(i)}}{\beta_a^{(i)}}.
\]

\section{Aggregation to a global Cobb--Douglas utility}

We seek a global utility for agent \(i\) of the form
\[
u_i(\bm{x}) = \prod_{a=1}^n x_a^{\beta_a^{(i)}},
\qquad \bm{x} \in \mathbb{R}_{++}^n,\quad \beta_a^{(i)} > 0,
\]
where
\[
\bm{\beta}^{(i)} \coloneqq (\beta_1^{(i)},\dots,\beta_n^{(i)})^\top \in \mathbb{R}_{++}^n
\]
is the (unnormalized) exponent vector. Normalizing to unit sum gives
\[
\tilde{\bm{\beta}}^{(i)} \coloneqq \frac{\bm{\beta}^{(i)}}{\bm{1}^\top \bm{\beta}^{(i)}},
\qquad \bm{1} \in \mathbb{R}^n,
\]
so that \(\sum_a \tilde{\beta}^{(i)}_a = 1\).

\subsection{Exact consistency}

The pairwise utilities are exactly consistent with some \(\bm{\beta}^{(i)}\) if and only if for all distinct \(a,b\)
\[
\alpha^{(i)}_{ab} = \frac{\beta_a^{(i)}}{\beta_a^{(i)}+\beta_b^{(i)}}
\quad\Longleftrightarrow\quad
\frac{\beta_b^{(i)}}{\beta_a^{(i)}} = r^{(i)}_{ab}.
\]
Equivalently, every cycle \(a_1 \to a_2 \to \dots \to a_k \to a_1\) must satisfy the cycle condition:
\[
\prod_{s=1}^{k} r^{(i)}_{a_s a_{s+1}} = 1,
\qquad (a_{k+1}=a_1).
\]
This is the multiplicative ``no-arbitrage'' condition in ratio space, ruling out intransitive preference cycles.

\subsection{Constructing \texorpdfstring{\(\bm{\beta}^{(i)}\)}{beta}}

If the ratios are consistent, \(\bm{\beta}^{(i)}\) is determined up to a positive scalar.
Fix a reference good (e.g.\ good \(1\)) and set \(\beta_1^{(i)} = 1\). Then for each \(a\),
choose any path \(1 \to \cdots \to a\) and propagate:
\[
\beta_a^{(i)} = \left(\prod_{\ell \text{ on path}} r^{(i)}_{\ell,\ell+1}\right)\beta_1^{(i)}.
\]
The cycle condition guarantees path-independence. Finally normalize to obtain \(\tilde{\bm{\beta}}^{(i)}\).

\subsection{Noisy or inconsistent pairwise inputs}

Empirical \(\alpha^{(i)}_{ab}\) will generally be noisy and violate the cycle condition. A convenient projection is to
work in log space. Let \(y_a^{(i)} \coloneqq \log \beta_a^{(i)}\) and
\(\rho^{(i)}_{ab} \coloneqq \log r^{(i)}_{ab}\). Then exact consistency is
\[
y_b^{(i)} - y_a^{(i)} = \rho^{(i)}_{ab}\quad \forall a\neq b,
\]
and a least-squares recovery solves
\[
\min_{\bm{y}^{(i)}\in\mathbb{R}^n}\;\sum_{1\le a<b\le n}\left((y_b^{(i)}-y_a^{(i)})-\rho^{(i)}_{ab}\right)^2,
\]
with a gauge fixing such as \(y_1^{(i)}=0\). Setting \(\beta_a^{(i)} \coloneqq e^{y_a^{(i)}}\) then yields a
best-fit exponent vector, which can be normalized.

\section{The \(n\)-good reaction--diffusion exchange economy}

\subsection{Reaction phase (endogenous transformation)}

Let \(\mathcal{R}\) be a finite set of reaction rules (technologies). Each rule \(r \in \mathcal{R}\) is specified by:
\begin{itemize}[leftmargin=*]
  \item an \emph{input} vector \(\bm{a}_r \in \mathbb{R}_+^n\),
  \item an \emph{output} vector \(\bm{b}_r \in \mathbb{R}_+^n\),
  \item an intensity (decision variable) \(x_{i,r}^t \ge 0\) for agent \(i\) at time \(t\).
\end{itemize}
Agent \(i\)'s post-reaction endowment is
\[
\bm{e}_i^{t,\mathrm{react}}
  = \bm{e}_i^{t}
    + \sum_{r \in \mathcal{R}} x_{i,r}^{t}\,(\bm{b}_r-\bm{a}_r),
\]
subject to feasibility (componentwise):
\[
\sum_{r \in \mathcal{R}} x_{i,r}^{t}\,\bm{a}_r \le \bm{e}_i^{t}.
\]
This defines a convex feasible set because the constraints are linear and the update is affine in \(\{x_{i,r}^t\}\).

\paragraph{Conservation remark.}
Reaction rules typically change the composition of \(\bm{E}\) over time. If one wishes to enforce conservation of a
linear invariant (e.g.\ mass/energy), introduce a matrix \(L\) such that \(L\bm{b}_r = L\bm{a}_r\) for all \(r\),
which guarantees \(L\sum_i \bm{e}_i^t\) is preserved.

\subsection{Diffusion phase (exchange)}

Let \(\bm{p}^t \in \mathbb{R}_{++}^n\) be a price vector at round \(t\).
Each agent chooses a \emph{per-good} trade fraction
\[
\bm{\lambda}_i^t = (\lambda_{i,1}^t,\dots,\lambda_{i,n}^t)^\top \in [0,1]^n,
\]
and posts a supply vector (Hadamard product \(\odot\)):
\[
\bm{s}_i^t = \bm{\lambda}_i^t \odot \bm{e}_i^{t,\mathrm{react}}.
\]
The retained inventory is \((\bm{1}-\bm{\lambda}_i^t)\odot\bm{e}_i^{t,\mathrm{react}}\).

Let \(\bm{z}_i^t \in \mathbb{R}_+^n\) denote the agent's acquired bundle from exchange. Define the trade budget
(wealth committed to trading) as
\[
w_i^t \coloneqq \bm{p}^t \cdot (\bm{\lambda}_i^t \odot \bm{e}_i^{t,\mathrm{react}}).
\]
The diffusion problem is
\begin{equation}
\label{eq:diffusion-problem}
\max_{\bm{\lambda}_i^t \in [0,1]^n,\;\bm{z}_i^t \ge 0}
\quad
u_i\!\left((\bm{1}-\bm{\lambda}_i^t)\odot\bm{e}_i^{t,\mathrm{react}} + \bm{z}_i^t\right)
\quad\text{s.t.}\quad
\bm{p}^t \cdot \bm{z}_i^t \le w_i^t.
\end{equation}

\subsection{Marshallian demand for fixed \texorpdfstring{\(\bm{\lambda}_i^t\)}{lambda}}

For fixed \(\bm{\lambda}_i^t\), the problem in \(\bm{z}_i^t\) is standard Cobb--Douglas with linear budget.
Let \(\tilde{\bm{\beta}}^{(i)}\) be the normalized exponents. The Marshallian demand is
\[
z_{i,a}^t(\bm{\lambda}_i^t)
= \frac{\tilde{\beta}_a^{(i)}}{p_a^t}\,w_i^t,
\qquad a=1,\dots,n.
\]
The post-trade holdings are
\[
x_{i,a}^t(\bm{\lambda}_i^t)
= (1-\lambda_{i,a}^t)e_{i,a}^{t,\mathrm{react}} + z_{i,a}^t(\bm{\lambda}_i^t).
\]

\subsection{Concave optimization over \texorpdfstring{\(\bm{\lambda}_i^t\)}{lambda}}

Define the log-utility as a function of \(\bm{\lambda}_i^t\):
\[
\ell_i(\bm{\lambda}_i^t)
\coloneqq \log u_i(\bm{x}_i^t(\bm{\lambda}_i^t))
= \sum_{a=1}^n \tilde{\beta}_a^{(i)}\log x_{i,a}^t(\bm{\lambda}_i^t).
\]
Since each \(x_{i,a}^t(\bm{\lambda})\) is affine in \(\bm{\lambda}\) and \(\log\) is concave, \(\ell_i\) is concave
on \([0,1]^n\). Hence the optimization in \(\bm{\lambda}_i^t\) is a convex program (maximize concave objective with
box constraints), solvable by projected gradient, coordinate ascent, or off-the-shelf convex solvers.

A useful closed-form expression for the partial derivatives is obtained by writing
\[
w_i^t = \sum_{b=1}^n p_b^t \lambda_{i,b}^t e_{i,b}^{t,\mathrm{react}},
\qquad
x_{i,a}^t = (1-\lambda_{i,a}^t)e_{i,a}^{t,\mathrm{react}} + \frac{\tilde{\beta}_a^{(i)}}{p_a^t}\,w_i^t.
\]
Then for each component \(m \in \{1,\dots,n\}\),
\[
\frac{\partial \ell_i}{\partial \lambda_{i,m}^t}
=
-\tilde{\beta}_m^{(i)}\frac{e_{i,m}^{t,\mathrm{react}}}{x_{i,m}^t}
+
p_m^t e_{i,m}^{t,\mathrm{react}}
\sum_{a=1}^n
\tilde{\beta}_a^{(i)}\,
\frac{\tilde{\beta}_a^{(i)}}{p_a^t\,x_{i,a}^t}.
\]
At an interior optimum, these partial derivatives vanish (KKT conditions), while boundary optima occur at
\(\lambda_{i,m}^t \in \{0,1\}\) for some components.

\section{Market clearing and price determination}

Total posted supply of good \(a\) is
\[
Q_a^t = \sum_{i=1}^m \lambda_{i,a}^t e_{i,a}^{t,\mathrm{react}}.
\]
Total demand (from Marshallian demand) is
\[
D_a^t = \sum_{i=1}^m \frac{\tilde{\beta}_a^{(i)}}{p_a^t}\,w_i^t,
\qquad
w_i^t=\bm{p}^t\cdot(\bm{\lambda}_i^t\odot\bm{e}_i^{t,\mathrm{react}}).
\]
Market clearing requires \(D_a^t = Q_a^t\) for all \(a\).

If all agents share the same exponent vector \(\tilde{\bm{\beta}}\) and choose a common scalar trading rate
\(\lambda_i^t\equiv\lambda^t\), then relative prices recover the familiar closed form
\[
\frac{p_a^t}{p_b^t}
=
\frac{\tilde{\beta}_a\,Q_b^t}{\tilde{\beta}_b\,Q_a^t}.
\]
In heterogeneous or fully vector-valued \(\bm{\lambda}_i^t\) settings, prices are naturally computed as a fixed point
(e.g.\ via tâtonnement or an excess-demand root-finder).

\section{Conclusion}

Aggregating pairwise Cobb--Douglas preferences into a global exponent vector enables a faithful \(n\)-good extension of
the two-good model in \parencite{Tristain2024}. The extended reaction--diffusion economy preserves:
\begin{itemize}[leftmargin=*]
\item convexity of feasible reaction sets and diffusion optimization,
\item closed-form Marshallian demand conditional on trading rates,
\item a principled interpretation of endogenous transformations (reaction) plus Pareto-improving exchange (diffusion),
\item and a clear pathway to simulation and implementation for \(n\ge 2\) goods.
\end{itemize}

\clearpage
\nocite{CobbDouglas1928,MasColell1995,BoydVandenberghe2004,Varian1992}

\printbibliography

\end{document}
